% ============================================================================
% PLANTILLA — WORKSHOP (taller práctico con software)
% Compilar:  scripts/build.sh workshop.tex   (XeLaTeX)
% ============================================================================
\documentclass{academic-beamer}
\title{Título del workshop}
\subtitle{Herramienta / software}
\author{Edison Achalma B.Sc. Econ.}
\institute{Universidad Nacional de San Cristóbal de Huamanga}
\date{\today}

\begin{document}
\begin{frame}\titlepage\end{frame}

\begin{frame}{Objetivos y requisitos}
\begin{itemize}\item Al terminar podrás\ldots \item Requisitos previos: software y datos\end{itemize}
\end{frame}

\begin{frame}[fragile]{Actividad 1 — Preparación}
\begin{codigo}[language=R]
library(tidyverse)
datos <- read_csv("datos.csv")
\end{codigo}
\begin{nota}Consigna: cargar los datos y describirlos.\end{nota}
\end{frame}

\begin{frame}[fragile]{Actividad 2 — Análisis}
\begin{codigo}[language=R]
modelo <- lm(y ~ x, data = datos); summary(modelo)
\end{codigo}
\end{frame}

\begin{frame}{Cierre}
\begin{itemize}\item Qué construimos \item Recursos para seguir practicando\end{itemize}
\end{frame}
\end{document}
